\chapter{Integração do motor Prolog com o componente Java}
\label{cap:integracao}

\cod{ExpertSystemBackend} define o ciclo comum aos motores: avançar o relógio,
iniciar o ciclo, inserir eventos, atualizar o veículo, decidir e expor a memória
e a explicação. \cod{PrologBackend} implementa esse contrato e permite trocar o
motor sem modificar o \cod{SimulationEngine}.

\paragraph{Comunicação.} O motor corre num processo \cod{swipl} persistente,
gerido pela classe \cod{PrologProcess}. O Java copia os ficheiros Prolog para
uma pasta de trabalho e arranca o processo com
\cod{swipl -q -g main -t halt server.pl}, que executa \cod{main/0} em \cod{server.pl}
e carrega os outros três ficheiros Prolog. O arranque acontece quando o motor é
iniciado, num RESET ou numa troca de motor de inferência. Os dois lados trocam mensagens
JSON por \cod{stdin} e \cod{stdout} padrão: um pedido por linha, e exatamente uma resposta
por linha, mesmo quando o pedido falha. Lançar o interpretador custa cerca de 150 a
250\,ms, o que é incompatível com uma decisão a cada 100\,ms. Além disso, o lado
Prolog tem de guardar estado entre ciclos. Por isso o processo é mantido vivo
durante toda a sessão. O pedido \cod{reset} inicia a sessão com a velocidade de
cruzeiro do cenário e com as constantes de paragem e as prioridades definidas no
Java, para que não haja valores duplicados nos dois lados. A
\cref{tab:pedidos-json} resume as mensagens JSON.

\begin{table}[H]
  \centering
  \caption{Pedidos JSON aceites pelo \texttt{server.pl}}
  \label{tab:pedidos-json}
  \small
  \begin{tabularx}{\textwidth}{@{}L{2.9cm}VL{5.0cm}VX@{}}
    \toprule
    \textbf{Comando} & \textbf{Entrada principal} & \textbf{Resposta} \\
    \midrule
    \cod{tick} & Veículo, eventos e decisão anterior
      & Propostas, factos, limite, verificações e métricas \\ \linhasep
    \cod{reset} & Cruzeiro, parâmetros e prioridades
      & Confirmação e reinicialização do estado \\ \linhasep
    \cod{whynot} & Identificador da regra
      & Condições que impediram a correspondência \\ \linhasep
    \cod{dependencyIndex} & Booleano
      & Estado do índice para \emph{benchmark} \\
    \bottomrule
  \end{tabularx}
\end{table}

O \cod{server.pl} lê uma linha, converte o JSON num \emph{dict} Prolog~\cite{swijson} e
responde sempre com outra linha. O \emph{stdout} contém apenas respostas.
Eventuais mensagens de aviso ou erro seguem para o \emph{stderr}. O \emph{flush}
explícito evita que o Java fique à espera de dados retidos no \emph{buffer}.

\paragraph{Pedidos e respostas.} Durante o ciclo, o \cod{PrologBackend} acumula
os eventos dos sensores e os comandos do condutor. Preenche também um
\cod{VehicleInput} com o que o Drools lê do facto do veículo: velocidade,
posição, faixa, estado da manobra e faixa de destino, marcações da estrada,
segurança das faixas adjacentes (calculada pela perceção em Java) e se há um
perigo à frente na faixa da direita. Tudo segue num único pedido \cod{tick}, com
a decisão anterior (ação, regra e prioridade). A resposta traz:

\begin{itemize}
  \item as propostas, cada uma com a ação, a velocidade, o risco, a prioridade,
    a razão, a regra e a evidência;
  \item os factos do ciclo;
  \item o limite de velocidade em vigor;
  \item os subobjetivos da segurança de mudança de faixa nos dois lados;
  \item as métricas do motor.
\end{itemize}

\paragraph{Arbitragem e aplicação.} O Java não decide nada por si: entrega as
propostas ao \cod{DecisionContext}. A
arbitragem trata em separado a velocidade e a faixa.

\begin{itemize}
  \item Na velocidade, as emergências e as paragens obrigatórias (prioridade 60
    ou mais) sobrepõem-se a tudo. Entre elas vence a mais prioritária e, em
    empate, a mais lenta. Abaixo disso, cada teto de velocidade que se verifica
    limita a velocidade de base, e decide o mais baixo, se ficar abaixo dessa
    base.
  \item Na faixa, vence a maior prioridade. Em empate, regressar à direita vence
    mudar de faixa.
  \item Uma mudança de faixa só é vetada quando uma paragem de prioridade
    superior está de facto a travar o veículo.
  \item Sem propostas, o veículo mantém a velocidade e a faixa.
\end{itemize}

Cada proposta fica com o seu resultado: venceu, perdeu, foi vetada ou foi
limitada. O \cod{DecisionActuator} aplica a decisão ao veículo. Pode recusar uma
mudança de faixa que não consiga executar, e essa recusa fica registada. O ciclo
é então guardado para o módulo de explicação.
